% !PRINTED 459
\weberpage{459}
%<<CONT
\[
24\,P_7,\ 56\,P_3,\ 84\,\bar P_2 \ \text{auf der Curve } f.
\]
\[
21\,P_8,\ 42\,P_4,\ 28\,P_6 \ \text{nicht auf der Curve } f.
\]
Alle anderen Punkte der Curve $f$ gehen durch die Substitutionen der
Gruppe in 168 verschiedene Punkte über.

Daraus geht noch unmittelbar hervor, dass die Curve $f$ keinen
Doppelpunkt haben und also um so weniger in Curven niedrigeren Grades
zerfallen kann. Denn angenommen, sie hätte einen Doppelpunkt, so
müssten auch alle seine Bildpunkte Doppelpunkte sein, und weil eine
Curve vierter Ordnung, auch wenn sie zerfällt, nicht mehr als sechs
Doppelpunkte haben kann, wenn sie nicht unendlich viele Doppelpunkte,
d.\,h. doppelt gezählte Curventheile hat, so müsste $f$ das Quadrat
einer quadratischen Form sein. Die Wurzel aus $f$ müsste dann auch
eine Invariante sein, während es doch keine quadratischen Invarianten
giebt.

Man kann übrigens auch leicht direct zeigen, dass die Grundcurve
keinen Doppelpunkt hat; denn die Bedingungen für einen Doppelpunkt
sind:
\[
f'(x_1) = 3x_1^2x_3 + x_2^3 = 0,
\]
\[
f'(x_2) = 3x_2^2x_1 + x_3^3 = 0,
\]
\[
f'(x_3) = 3x_3^2x_2 + x_1^3 = 0,
\]
und diese Gleichungen können nicht anders erfüllt sein, als wenn
$x_1, x_2, x_3$ verschwinden (\S.\,86).

Der Kürze wegen wollen wir ein System von Punkten, deren jeder in
jeden anderen durch Substitutionen der Gruppe transformirbar ist, ein
\gesperrt{System verbundener Punkte} nennen.

\S.\,123.

\gesperrt{Die höheren Invarianten.}

Weitere Invarianten unserer Gruppe lassen sich nach dem Satze 4.,
\S.\,40 ableiten, indem wir Covarianten der Form $f$ bilden. Wir
fassen die in \S.\,87 allgemein besprochenen Covarianten ins Auge, und
bilden zunächst die \gesperrt{Hesse}'sche Covariante
\weq{1}{\Delta = \frac{1}{54}
\begin{vmatrix}
f''(x_1,x_1), & f''(x_1,x_2), & f''(x_1,x_3) \\
f''(x_2,x_1), & f''(x_2,x_2), & f''(x_2,x_3) \\
f''(x_3,x_1), & f''(x_3,x_2), & f''(x_3,x_3)
\end{vmatrix},}
%CONT>>
